\subsection{An exact comparison for the unrestricted alternating-path model}

In E.20--E.24 we take $\Sigma=\Mor(X)$. We give the comparison for the precise paths and maps under discussion, rather than replace them silently by another categorical cocylinder. Let $J$ be the integer fence: its objects are $j\in\Z$, its nonidentity arrows are
\[
                 2k\longrightarrow 2k-1,
                 \qquad 2k\longrightarrow 2k+1,
\]
and there are no composable pairs of nonidentity arrows. Put $\mathcal P(X)\subseteq\Fun(J,X)$ for the full subcategory of functors which are constant, with identity arrows, outside some finite interval. A natural transformation between two such functors is necessarily constant on each sufficiently distant tail. Its two eventual components define the endpoint functor
\[
 p=(\sigma,\beta):\mathcal P(X)\longrightarrow X\times X.
\]
Thus $\mathcal P(X)$ is the manuscript's two-sided $\Ch_\infty(X)$ with this explicit parity convention. Its strict fibre $\mathcal P(X;a,b)$ has transformations whose two eventual components are $1_a,1_b$. We write $B$ for classifying space and $K=BX$. All homotopy categories of spaces over a base below mean the localization of the \emph{slice category} at underlying weak equivalences, not the slice of the ordinary homotopy category.

\begin{claim}{E.20}{Proved}{Finite support and the constant-path equivalence}
The constant-path section $j:X\to\mathcal P(X)$ is a weak equivalence. For each $a\in X$, both one-endpoint categories $\mathcal P(X;a,?)$ and $\mathcal P(X;?,a)$ are weakly contractible.
\end{claim}
\begin{proof}
For integers $n\leq m$, let $\mathcal P_{n,m}(X)$ be the full subcategory of paths constant outside $[n,m]$. Restriction identifies it with $\Fun(J|_{[n,m]},X)$. Extension by identities gives fully faithful transition functors when the interval increases. These categories form a filtered union on objects and arrows; every finite diagram is contained in one stage. The same statements hold with either or both endpoints fixed.

A finite oriented interval $I$ with vertices $0,\ldots,r$ has functors $r_k:I\to I$, $r_k(i)=\min(i,k)$, for $0\leq k\leq r$. They preserve each directed generating edge, sending an edge beyond $k$ to an identity. The successive functors $r_k,r_{k+1}$ are joined by a natural transformation in the direction of the edge between $k$ and $k+1$: its components on the tail are that edge, and its other components are identities. On every generating edge the naturality square commutes. Precomposition gives a finite zigzag from the constant-at-initial-value endofunctor of $\Fun(I,X)$ to its identity. This zigzag fixes the initial endpoint. Consequently the constant-diagram functor $X\to\Fun(I,X)$ is a weak equivalence; with initial endpoint fixed at $a$ the functor category is contractible. Reversing the enumeration proves the version with the terminal endpoint fixed.

Nerves preserve filtered colimits of categories, since the data and equations of a composable finite string occur at a stage. Weak equivalences of simplicial sets are preserved by filtered colimits. For completeness, $\operatorname{Ex}^{\infty}$ commutes with filtered colimits: its individual simplex formulas use finite subdivided simplices, and its sequential iteration also commutes with such colimits. In the resulting Kan complexes, homotopy groups and components are detected by finitely many simplices and finite homotopies. Applying this observation to the stagewise constant-diagram equivalences proves the first assertion and the two assertions with one endpoint fixed. No contraction of a strict two-endpoint fibre has been claimed.
\end{proof}

\begin{claim}{E.21}{Proved}{The exact endpoint fibres are homotopy fibres}
For each $a,b\in X$, the natural square
\[
\begin{array}{ccc}
 B\mathcal P(X;a,b)&\longrightarrow&B\mathcal P(X)\\
 \downarrow&&\downarrow Bp\\
 \{(a,b)\}&\longrightarrow&K\times K
\end{array}
\]
is homotopy Cartesian. In particular $B\mathcal P(X;a,b)$ has the weak homotopy type of the ordinary path space $P_{a,b}K$. The comparison with the homotopy fibre is natural for functors over $X\times X$.
\end{claim}
\begin{proof}
We first establish the two lifting properties for this exact category. If $n<m$ are even, both endpoints of $J|_{[n,m]}$ are sources. The endpoint functor on $\Fun(J|_{[n,m]},X)$ is a Grothendieck fibration: to lift $(u,v):(a',b')\to(a,b)$ into a diagram with endpoints $(a,b)$, replace the endpoint objects by $a',b'$ and precompose its two outer arrows by $u,v$. Leave every interior object and arrow unchanged. The transformation to the original diagram has components $u,v$ at the endpoints and identities inside. Any transformation over $(u,v)$ factors uniquely through it, so it is Cartesian. Positive length is essential here. If $n<m$ are odd, both endpoints are sinks, and postcomposition supplies the dual coCartesian construction.

For a general functor $q:E\to B_0$, an object $e\in E$, and an arrow $h:b\to q(e)$, write $L_q(e,h)$ for the category of pairs $(d,\gamma)$ with $q(d)=b$, $q(\gamma)=h$, and $\gamma:d\to e$; its arrows are vertical maps commuting with $\gamma$. For the functor $p$ above, this is the filtered union of the lift categories at sufficiently large positive even-even support intervals. Each has a terminal Cartesian lift, and hence is contractible. Every object, arrow, and equation in a finite diagram of lifts occurs at one such stage. The filtered-colimit argument of E.20 shows that $L_p(e,h)$ is weakly contractible. Odd-odd intervals prove, dually, that every outgoing-lift category is weakly contractible. Both parity families are cofinal among all finite supports. This proves the usual smoothness and properness lift tests directly; it does not assert that the union itself is a Grothendieck fibration or opfibration.

We spell out why these lift tests imply the required Quillen--B condition. Let $q:E\to B_0$ have both properties just proved, and let $E_b$ be its strict fibre. The inclusion
\[
                 i_b:E_b\longrightarrow b\backslash q
\]
is a weak equivalence by Quillen A: its over-comma category at $(e,h:b\to q(e))$ is exactly $L_q(e,h)$. For an arrow $h:b\to b'$, let $M_h$ have as objects arrows $\gamma:d\to d'$ of $E$ over $h$, and as arrows commuting squares whose two vertical sides lie over identities. Let $s:M_h\to E_b$ and $t:M_h\to E_{b'}$ be the source and target projections.

For $d_0\in E_b$, an object of $d_0\backslash s$ is an arrow $\gamma:d\to d'$ over $h$, together with a vertical $\alpha:d_0\to d$. Composition sends it to the outgoing lift $\gamma\alpha:d_0\to d'$. The inclusion of that outgoing-lift category takes $\alpha=1_{d_0}$. The composite in one direction is the identity; in the other, the square $(\alpha,1_{d'})$ gives a natural transformation from the composite to the identity. Thus $d_0\backslash s$ is homotopy equivalent to the contractible outgoing-lift category. Quillen A makes $s$ a weak equivalence. An object of $t/d'_0$ instead consists of $\gamma:d\to d'$ and a vertical $\alpha:d'\to d'_0$. The functor $(\gamma,\alpha)\mapsto\alpha\gamma$ retracts this category onto the incoming lifts into $d'_0$. Its inclusion takes $\alpha=1_{d'_0}$, and the squares $(1_d,\alpha)$ give the natural transformation from the identity to the other composite. Thus $t/d'_0$ is contractible and Quillen A makes $t$ a weak equivalence.

The arrows $\gamma$ themselves give a natural transformation
\[
             i_b s\ \Longrightarrow\ h^*i_{b'}t
             \quad\hbox{into }b\backslash q,
\]
where $h^*:b'\backslash q\to b\backslash q$ is precomposition. As $i_b,s,i_{b'},t$ are weak equivalences, this natural transformation and two-out-of-three show that $h^*$ is a weak equivalence. Quillen B therefore identifies $B(b\backslash q)$ with the homotopy fibre of $Bq$ over $b$, and the natural map from $BE_b$ to that homotopy fibre is a weak equivalence \citep[\S1, Theorems A and B]{Quillen1973}. This formulation is also recorded in \citet[Theorem 1.8]{Shoikhet2016}. The comma constructions and the maps to a standard homotopy-fibre model commute with functors over the base, which gives the asserted naturality.

Apply this argument to $p$. E.20 identifies the constant-path section as a weak equivalence over $K^2$, from the diagonal $K\to K^2$ to $Bp$. Consequently their derived fibres agree. The homotopy fibre of the diagonal over $(a,b)$ is the path space from $a$ to $b$. Equivalently, let $P(K)=K^{[0,1]}$ with its endpoint fibration and constant-path section $s:K\to P(K)$. The comparison is the isomorphism
\[
 u_{\mathcal P}=s\,(Bj)^{-1}:B\mathcal P(X)\longrightarrow P(K)
       \quad\hbox{in }\operatorname{Ho}(\mathrm{Spc}/K^2),
\]
followed by the derived-fibre functor. The preceding proof identifies its source derived fibre with the stated strict categorical fibre.
\end{proof}

\begin{claim}{E.22}{Proved}{The actual even translation is the identity in the homotopy category}
Let $T^2:\mathcal P(X)\to\mathcal P(X)$ be given by $(T^2c)(j)=c(j-2)$, with the corresponding shift on arrows and on natural transformations. It induces the identity in $\Hot$ on the total category and on every fixed-endpoint category.
\end{claim}
\begin{proof}
Translation by two is an automorphism of the integer fence, preserves eventual endpoint values, and fixes every constant diagram and constant natural transformation. Thus $T^2$ is an exact endofunctor over $X^2$ satisfying $T^2j=j$. The morphism $Bj$ is invertible in the localized slice $\operatorname{Ho}(\mathrm{Spc}/K^2)$ by E.20. Cancellation gives $B(T^2)=1$ there. Apply the derived-fibre functor at $(a,b)$. Naturality in E.21 identifies the resulting map with the classifying-space map of the \emph{actual restriction} of $T^2$ to $\mathcal P(X;a,b)$. Hence that map is the identity in $\Hot$. The argument neither assumes faithfulness of finite tests nor constructs a bounded zigzag of natural transformations from the two functors.
\end{proof}

\begin{claim}{E.23}{Proved}{The one-sided models and their actual comparison maps}
Put $D=(+,-)$. Let $\mathcal L(X)$ be the colimit of the categories of $D^n$-paths under padding by a constant $D$-block on the left, and let $\mathcal R(X)$ use padding on the right. They have constant-path sections and endpoint maps. Their constant-path sections are weak equivalences, their one-endpoint fibres are weakly contractible, and their strict two-endpoint fibres compute their derived endpoint fibres. Constant-tail extension gives functors
\[
 q_L:\mathcal L(X)\longrightarrow\mathcal P(X),
 \qquad q_R:\mathcal R(X)\longrightarrow\mathcal P(X)
\]
which are weak equivalences on total categories and on every fixed-endpoint fibre. These comparisons commute with the canonical path-space comparisons of E.21.
\end{claim}
\begin{proof}
Realize $\mathcal L(X)$ on the half-fence $j\leq0$, eventually constant at $-\infty$, using stages $[-2n,0]$; realize $\mathcal R(X)$ on $j\geq0$, using $[0,2n]$. Thus $q_L,q_R$ literally extend the other half by the value at zero. Padding is fully faithful, and all the constant-section and one-endpoint assertions follow from the interval contractions and filtered-colimit proof of E.20.

We prove the comparison on the actual half-fence by sliding an identity collar. Let $S:\mathcal L(X)\to\mathcal L(X)$ append two identity arrows at the right end, shifting the old path two places left. Thus $(Sc)(i)=c(\min(i+2,0))$. We give a homotopy $BS\simeq1$, relative to both endpoints. For $j\geq0$, define an endofunctor $\phi_j$ of the negative half-fence by
\[
 \phi_j(i)=
 \begin{cases}
 i+2,&i\leq-j-2,\\
 -j,&-j-2\leq i\leq-j,\\
 i,&i\geq-j.
 \end{cases}
\]
The values agree at the boundaries. Each generating arrow is either translated by two, unchanged, or sent to an identity; thus $\phi_j$ is a functor and fixes zero. Precomposition gives $F_j(c)=c\phi_j$. We have $F_0=S$. For even $j$ there is a natural transformation $F_j\to F_{j+1}$; for odd $j$ its direction is reversed. Its only possibly nonidentity components are at $-j-2,-j-1$, where both are the edge between $c(-j)$ and $c(-j-1)$. The naturality squares on the three adjacent generating edges commute; all other squares have identity components. The two endpoint components are identities.

For every finite diagram in $\mathcal L(X)$, the functors $F_j$ and these transformations become literally the identity on that diagram for sufficiently large $j$, because its paths and transformations are constant sufficiently far to the left. Concatenate the nerve homotopies on successive time intervals with endpoints $1-2^{-j}$, traversing a natural-transformation homotopy backwards when necessary, and put the identity at time $1$. This is continuous: on each nerve simplex it is stationary after finitely many intervals. In compactly generated spaces, realization has the weak topology with respect to its cells, and product with the compact interval preserves this cellwise test. We obtain $BS\simeq1$, on the total category and on every fixed-endpoint category. A related finite-stage identity-pair sliding calculation appears in \citet[\S3.2.1, Proposition 3.3]{Shoikhet2016}; the explicit formulas above establish the required homotopy for our colimit model independently of a categorical cocylinder identification.

Now $q_LS=T^{-2}q_L$ strictly. Every finite diagram in $\mathcal P(X)$ is carried into the image of $q_L$ by a sufficiently large negative even translation. The same holds for any finite diagram representing a homotopy. Homotopy groups and components of these nerves are computed at finite support by the filtered simplicial argument of E.20. These facts prove that $q_L$ is a weak equivalence, including on every fixed-endpoint fibre. Here is the basepoint detail. For a vertex $c\in\mathcal L(X;a,b)$, the collar homotopy supplies paths from $c$ to $S^kc$ inside that category. Use those paths to identify the relevant based homotopy groups. Translation $T^{-2k}$ gives an automorphism of the homotopy group of $B\mathcal P(X;a,b)$, and the relation $q_LS^k=T^{-2k}q_L$ shows that it preserves the image of the homotopy group of $B\mathcal L(X;a,b)$, in both directions, because $S^k$ is a homotopy equivalence. Any class translates into that image at some finite support, so it was already in the image. If a class from the half-fence becomes zero, translate a finite nullhomotopy into the half-fence; it kills its image under $S^k$, and invertibility of $S^k$ kills the original class. The same argument with paths proves bijectivity on components. Reversal of the fence proves the result for $q_R$.

The total and fibrewise weak equivalences $q_L,q_R$, together with E.21 and naturality of the maps to homotopy fibres, transfer the strict-fibre theorem to $\mathcal L$ and $\mathcal R$. Explicitly, the comparison square has the strict-fibre map on top, the map of homotopy fibres induced by $Bq_L$ or $Bq_R$ below, and the natural fibre comparisons vertically. Its top and bottom maps and its target vertical map are weak equivalences, so its source vertical map is one as well. This argument does not assert that a one-sided endpoint map has both strict categorical lifting properties.

If $j_L,j_R$ denote the constant sections, $q_Lj_L=j$ and $q_Rj_R=j$. Cancellation in the localized slice therefore gives $u_{\mathcal P}Bq_L=s(Bj_L)^{-1}$ and likewise for $q_R$, proving compatibility for these actual maps.

We will also use the elementary iterated-fibre identity. For $(f_1,f_2):A\to B_1\times B_2$, if $A_b=f_1^{-1}(b)\to\operatorname{hofib}_b(f_1)$ is a weak equivalence, regard both spaces over $B_2$ by $f_2$. Their derived fibres over $c$ agree; the latter is $\operatorname{hofib}_{(b,c)}(f_1,f_2)$. This follows from the model consisting of a point of $A$ and a pair of paths in the two base spaces. The maps from strict fibres in this identity are natural.
\end{proof}

\begin{claim}{E.24}{Proved}{The manuscript's left-right concatenation is associative}
For the block $D=(+,-)$ and $\Sigma=\Mor(X)$, glue a left-half path and a right-half path at their common value at zero. This is an actual functor
\[
 \kappa:\mathcal L(X)\times_X\mathcal R(X)\longrightarrow\mathcal P(X),
\]
where the pullback uses the terminal and initial endpoints. On fixed fibres define the manuscript's composition by
\[
 \mu_{a,b,c}=B\kappa\,\bigl((Bq_L)^{-1}\times(Bq_R)^{-1}\bigr):
 B\mathcal P(X;a,b)\times B\mathcal P(X;b,c)
       \longrightarrow B\mathcal P(X;a,c)
 \quad\hbox{in }\Hot.
\]
Here the inverses refer to the fixed-fibre weak equivalences of E.23. Under the canonical comparisons to ordinary path spaces, this \emph{particular} operation is ordinary path concatenation. It is therefore associative and unital in $\Hot$; reversal supplies inverses in the usual path-groupoid sense, and the loop object is a group object of $\Hot$.
\end{claim}
\begin{proof}
The gluing definition agrees on the common vertex and on the common component of a natural transformation, so it defines a functor on the indicated strict pullback. At stages $[-2n,0]$ and $[0,2n]$ it is exactly concatenation of two $D^n$-paths into a $D^{2n}$-path. Increasing $n$ pads at the two outer ends, exactly as in the manuscript's formula. This identifies the specified map, including its transition maps.

Write $E=\mathcal L(X)\times_X\mathcal R(X)$ and let $r:E\to X^3$ record the initial, middle, and terminal values. Its constant-pair section $j_E:X\to E$ is a weak equivalence. One proof uses the filtered system of finite supports: a stage with left length $2n$ and right length $2m$ is the functor category on their glued oriented interval, and its constant section is a weak equivalence by the same explicit interval contraction as in E.20. The constants commute with all enlargements.

We need a relative assertion as well. The middle-coordinate map $BE\to K$ is a weak equivalence since its composite with $Bj_E$ is $1_K$. Its strict fibre at $b$ is
\[
 B\mathcal L(X;?,b)\times B\mathcal R(X;b,?),
\]
which is contractible by E.23. Hence this strict fibre computes the middle-coordinate homotopy fibre. On it the remaining two coordinates are the product of the initial-endpoint map on $\mathcal L(X;?,b)$ and the terminal-endpoint map on $\mathcal R(X;b,?)$. Their strict fibres compute homotopy fibres: apply the iterated-fibre identity in E.23 to each full endpoint map, whose strict-fibre assertion is already proved, using its contractible one-endpoint fibre. Taking the remaining two derived fibres, and using that homotopy fibres commute with products, proves
\[
 B\bigl(\mathcal L(X;a,b)\times\mathcal R(X;b,c)\bigr)
       \simeq \operatorname{hofib}_{(a,b,c)}(BE\xrightarrow{Br}K^3).
\]
Every comparison in this paragraph is the natural comparison from the strict fibre, so it also controls the maps induced by projections and gluing.

Let $P=K^{[0,1]}$ and $Q=P\times_K P$, with their endpoint maps to $K^2$ and $K^3$. Both endpoint maps are fibrations; $Q$ is the space of maps from $[0,2]$ with values specified at $0,1,2$. Write $s:K\to P$ and $t:K\to Q$ for constants. These are weak equivalences. Define in the localized slices
\[
 u_L=s(Bj_L)^{-1},\qquad u_R=s(Bj_R)^{-1},\qquad
 u_{\mathcal P}=s(Bj)^{-1},\qquad v=t(Bj_E)^{-1}.
\]
For each projection from $E$ and for $\kappa$, regard the target over $K^3$ by adjoining the unused endpoint coordinate. This amounts to a product with $K$, which preserves weak equivalences. The two maps comparing categorical gluing with topological concatenation agree after composition with $Bj_E$: each sends a constant pair to a constant path, retaining the same middle coordinate. Since $Bj_E$ is invertible in this slice, cancellation proves
\[
                  u_{\mathcal P}B\kappa
                       =\operatorname{concat}\,v.
\]
The same cancellation proves that $v$ commutes with both projections under $u_L,u_R$.

Apply the derived-fibre functor at $(a,b,c)$. The relative assertion for $E$, E.21, E.23, and the ordinary path fibrations identify the derived fibres with the stated strict fibres. The two projection identities identify the induced map of the product fibre as $u_{L,a,b}\times u_{R,b,c}$, because maps into a product in $\Hot$ are determined by their two projections. Thus
\[
 u_{\mathcal P,a,c}\,B\kappa
 =\operatorname{concat}\,(u_{L,a,b}\times u_{R,b,c}).
\]
Compatibility with $q_L,q_R$ from E.23 gives the claimed formula for $\mu$. Ordinary concatenation is associative up to the endpoint-fixed homotopy which changes the two subdivision points in a threefold concatenation. Constants are units up to the analogous reparameterization homotopy. The reversal $c(j)\mapsto c(-j)$ of the integer fence is an actual functor, swaps endpoints, and preserves the constant section; cancellation in the same slices identifies it with reversal of topological paths. The standard contraction of a path followed by its reverse supplies the inverse identities. These prove all the asserted identities in $\Hot$ for the specified operation.
\end{proof}

\subsection{Relative collars and the remaining quantified comparisons}

We now make explicit the greater generality of VII ms.~108--136. An \emph{admissible} class $\Sigma\subseteq\Mor(X)$ contains every identity and is closed under composition. For a finite sign word $\tau$, let $I_\tau$ be the free category on its oriented interval. A $\Sigma$-path of type $\tau$ is a functor $I_\tau\to X$ taking every geometrically backward arrow to $\Sigma$. The category of such paths is the \emph{full} subcategory of $\Fun(I_\tau,X)$ on those objects. In particular, an interior component of a natural transformation is not required to belong to $\Sigma$; fixed-endpoint transformations have identity components at the two ends. These are the manuscript's definitions at VII ms.~65--67.

\begin{claim}{E.25}{Proved}{Canonical relative collars and literal even translation}
For any admissible $\Sigma$, let $\mathcal L_\Sigma(X;a,b)$, $\mathcal R_\Sigma(X;a,b)$, and $\mathcal P_\Sigma(X;a,b)$ be the full relative subcategories of the canonical $D=(+,-)$ models of E.23. The actual constant-tail extension maps $q_L,q_R$ are weak equivalences on these fibres and on their total categories. The literal translation $(T^2c)(j)=c(j-2)$ induces the identity in $\Hot$ on every such fibre. Gluing with the inverses of $q_L,q_R$ consequently defines the canonical relative operation $\mu_D$ in $\Hot$; constant paths are its left and right units.
\end{claim}
\begin{proof}
Each $\phi_j$ in E.23 is nondecreasing on the geometric vertex order and sends a generating edge to an edge of the same geometric direction or to an identity. Hence $F_j$ preserves relative paths. The homotopies between successive $F_j$ are also permitted, since the relative path category is full and their endpoint components are identities. The eventual-stationarity argument therefore proves $BS\simeq1$ on each relative left-half-fence fibre.

There is an equally explicit right-hand version, without invoking path reversal, which need not preserve a relative path class. For $i,j\geq0$ put
\[
 \psi_j(i)=
 \begin{cases}
 i,&i\leq j,\\
 j,&j\leq i\leq j+2,\\
 i-2,&i\geq j+2.
 \end{cases}
\]
Precomposition gives $G_j$. The functor $G_0=A$ inserts two identities at the left end of a right-half path. Consecutive $G_j$ have their two possibly nonidentity comparison components at $j+1,j+2$, given by the edge between $c(j)$ and $c(j+1)$; the transformation points $G_j\to G_{j+1}$ for even $j$ and backwards for odd $j$. These functors preserve relative paths for the same direction-preservation reason. They become the identity on each finite diagram sufficiently far to the right. Thus $BA\simeq1$, relative to the endpoints.

Even translation preserves the relative class and is an automorphism of the two-sided model. The finite-support and basepoint argument of E.23 now applies word for word, with $q_LS=T^{-2}q_L$ and $q_RA=T^2q_R$, to prove both comparison functors weak equivalences. This use of E.23 requires no endpoint-fibration theorem. Cancelling $Bq_L$ in $Bq_L\,BS=B(T^{-2})Bq_L$ gives $B(T^{-2})=1$, hence $B(T^2)=1$, on the actual relative fibre.

The gluing functor preserves the relative condition, so $\mu_D$ is well defined using the two inverses in $\Hot$. Gluing a left-half path to the constant right-half path is exactly $q_L$; gluing a constant left-half path to a right-half path is exactly $q_R$. Both comparisons preserve constants. Cancellation proves the two unit identities. We have not used the derived-fibre argument of E.24 to claim relative associativity: its strict endpoint-fibre hypothesis was established there only for the all-arrows case.
\end{proof}

To specify the remaining questions, we use the actual type category $\mathrm{Ty}$. Its objects are finite sign words, including the empty word $e$. A morphism is an endpoint-preserving nondecreasing functor between the associated oriented interval categories. Equivalently, if the source word is $\varepsilon_1\cdots\varepsilon_r$, a morphism is a choice of integers $l_i\geq0$ whose target word is $\varepsilon_1^{l_1}\cdots\varepsilon_r^{l_r}$; the $i$th source edge maps to that same-sign string, or to an identity when $l_i=0$. This is the description at VII ms.~50'--53. The \emph{useful length} $\ell_u(\tau)$ is the number of maximal constant-sign runs in the word, with $\ell_u(e)=0$ (VII ms.~75--76).

A projective sequence $t:\N^{\op}\to\mathrm{Ty}$ is \emph{left-cofinal in the manuscript's sense} when
\[
       \forall\tau\in\mathrm{Ty}\quad
       \exists n\quad\Hom_{\mathrm{Ty}}(t_n,\tau)\ne\varnothing.
\]
This is cofinality for the reachability preorder of VII ms.~99--102, not an assumption that comma categories are contractible. Equivalently $\ell_u(t_n)$ is unbounded, as stated at VII ms.~138: a source with sufficiently many alternating sign runs maps to any prescribed word by collapsing unused runs and lengthening the retained same-sign runs; a morphism cannot increase useful length. For such $t$ define the actual category
\[
 C_{\infty,\Sigma}(t;a,b)
   =\colim_n\Fun_\Sigma(I_{t_n},X)_{a,b},
\]
with transition functors given by precomposition with the specified transitions of $t$.

\begin{claim}{E.26}{Open}{The remaining cofinal-type and all-type comparisons}
For arbitrary admissible $\Sigma$ and arbitrary left-cofinal projective type sequences as just defined, this audit leaves undecided the following \emph{specified geometric} assertions of VII ms.~108--123:
\begin{enumerate}
\item A natural transformation $u:t\to t'$ of projective sequences, consisting of strictly commuting maps $u_n:t_n\to t'_n$, always induces a weak equivalence
\[
 u^*:C_{\infty,\Sigma}(t';a,b)\longrightarrow
                  C_{\infty,\Sigma}(t;a,b).
\]
\item The isomorphisms obtained from the canonical type projections $t*t'\to t,t'$ are independent of all comparison choices and form the transitive canonical system asserted at ms.~119--123.
\item The manuscript's comparison from a sequential colimit to its \emph{all-type} path category is an isomorphism in $\Hot$ for every such relative datum. Precisely, let $E_t$ be the fibred Grothendieck construction of the displayed sequence over $\N^{\op}$ and let $\mathrm{Ch}_\Sigma(X;a,b)\to\mathrm{Ty}$ have fibre $\Fun_\Sigma(I_\tau,X)_{a,b}$. The assertion is that
\[
 E_t\longrightarrow\mathrm{Ch}_\Sigma(X;a,b)
\]
is a weak equivalence, so that the span with the ordinary-colimit comparison $E_t\to C_{\infty,\Sigma}(t;a,b)$ yields the claimed isomorphism.
\end{enumerate}
\end{claim}
\begin{proof}[Status and missing hypotheses]
The ordinary-colimit leg of the last span is covered by E.8; its fibred convention is obtained from the covariant one by taking opposite categories. That does not prove its all-type projection is a weak equivalence. The comparison maps in the first two items exist and are the stated precomposition functors, but an isomorphism of their ind-objects does not establish their underlying weak equivalence or equality in $\Hot$; E.4 and E.8 explain the general failure of that inference. E.23 and E.25 establish the concrete collar comparisons, not this entire quantified family of transition systems. No counterexample to these particular geometric assertions, and no complete independent proof of their asserted generality, is supplied here. Thus ``Open'' records the precise remaining audit question, not a claim that the literature has designated it as an open problem.
\end{proof}

\begin{claim}{E.27}{Open}{The remaining relative and arbitrary-block composition assertions}
Fix an admissible $\Sigma$ and a word $\tau_0$ with $\ell_u(\tau_0)\geq2$. Form the three actual systems of VII ms.~125--126: the stages $\tau_0^n$ with transitions which collapse only the extra left blocks, the stages $\tau_0^n$ with transitions which collapse only the extra right blocks, and the stages $\tau_0^{2n}$ with symmetric outer collapse. Denote their colimits by $L_{\tau_0,\Sigma}$, $R_{\tau_0,\Sigma}$, and $P_{\tau_0,\Sigma}$, with endpoints understood. Constant padding defines actual maps $q_L:L_{\tau_0,\Sigma}\to P_{\tau_0,\Sigma}$ and $q_R:R_{\tau_0,\Sigma}\to P_{\tau_0,\Sigma}$, and concatenation defines the actual gluing functor $\kappa_{\tau_0}$. Beyond E.24--E.25, the audit does not settle all of the following claims of ms.~125--127: both padding maps are weak equivalences for every block; the resulting operations $\kappa_{\tau_0}(q_L^{-1}\times q_R^{-1})$ are independent of the block under the specified comparison maps; and those operations are associative in $\Hot$ for every admissible $\Sigma$.
\end{claim}
\begin{proof}[Status and smallest uncovered extensions]
For $\tau_0=D=(+,-)$, E.25 proves that the relative operation exists and is unital. The additional assertion of its associativity for an arbitrary proper admissible class $\Sigma\subsetneq\Mor(X)$ is already an extension not proved here. The proof of E.24 cannot simply be reused: it requires the actual strict relative endpoint fibres to be derived fibres of the diagonal on $BX$, a hypothesis absent for general $\Sigma$. E.9 supplies a different derived mapping-space construction only under its explicit rectification hypothesis and does not identify this raw operation with it.

Even with $\Sigma=\Mor(X)$, the written canonical-block argument does not verify every specified comparison for a different block. An explicit short uncovered convention is the alternating word $(+,-,+)$, whose useful length is three, with the left, right, and symmetric transitions specified in the claim. This is a smallest extension of block length beyond the two-letter canonical case, not a claim of a counterexample. A proof of its comparison and composition compatibility, or a proof for the full quantified family, would settle that extension. The general independence and associativity assertions cannot be inferred from their images under $\psi$. We retain them here with their original quantifiers so that the established all-arrows canonical result does not silently replace them.
\end{proof}
